\chapter{Frontend}

Il frontend di Ruggine è sviluppato in React per garantire una UI moderna, reattiva e facilmente estendibile. La scelta di React è motivata dalla sua ampia diffusione, dalla ricchezza di librerie e dalla possibilità di strutturare l'applicazione in componenti riutilizzabili e testabili.

\section{Struttura interna}
La struttura del frontend è suddivisa in cartelle tematiche. La cartella \texttt{components} contiene i componenti riutilizzabili dell'interfaccia, come LoginForm, MessagesContainer, InputArea, Navbar e altri. La cartella \texttt{pages} raccoglie le pagine principali dell'applicazione, come LoginPage, ChatPage e InvitesPage, ognuna composta da più componenti. La cartella \texttt{context} gestisce lo stato globale tramite React Context, in particolare la connessione WebSocket e l'autenticazione (WebSocketContext.jsx). Infine, la cartella \texttt{utils} è destinata a funzioni di utilità e helper.

\subsection{Esempio: LoginForm}
	exttt{LoginForm.jsx} gestisce l'input delle credenziali e l'invio tramite callback. La separazione tra logica di presentazione e gestione dello stato permette di riutilizzare il componente sia per login che per registrazione.

\begin{lstlisting}[language=JavaScript, caption={Estratto da LoginForm.jsx}]
function LoginForm({ isLoginMode, username, setUsername, password, setPassword, handleApiCall }) {
	return (
		<form onSubmit={handleApiCall}>
			<input ... />
			<input ... />
			<button type="submit">{isLoginMode ? 'Accedi' : 'Registrati'}</button>
		</form>
	);
}
\end{lstlisting}

\section{Comunicazione e flussi}
La comunicazione con il backend avviene principalmente tramite API HTTP per operazioni come login, registrazione e invio inviti, e tramite WebSocket per la chat in tempo reale e gli aggiornamenti di gruppo/inviti. Il context globale (WebSocketContext.jsx) incapsula la logica di connessione, gestione messaggi, inviti e autenticazione, esponendo API semplici ai componenti figli.

\subsection{Esempio: invio messaggio}
\begin{lstlisting}[language=JavaScript, caption={Invio messaggio tramite context}]
const { sendMessage } = useWebSocket();
sendMessage(groupId, "Ciao a tutti!");
\end{lstlisting}

\section{Gestione dello stato e UI}
Lo stato globale (utente, gruppi, messaggi, inviti) è gestito tramite React Context e hook personalizzati. Questo approccio evita il "prop drilling" e rende la UI reattiva ai cambiamenti di stato.

\subsection{Esempio: visualizzazione messaggi}
\texttt{MessagesContainer.jsx} mostra i messaggi di una chat, distinguendo tra messaggi inviati e ricevuti:
\begin{lstlisting}[language=JavaScript, caption={Estratto da MessagesContainer.jsx}]
function MessagesContainer({ activeMessages, user, messagesEndRef }) {
	return (
		<div>
			{activeMessages.map((msg, idx) => (
				<div className={msg.from === user ? 'sent' : 'received'}>{msg.content}</div>
			))}
		</div>
	);
}
\end{lstlisting}

\section{Sicurezza e validazione}
\begin{itemize}
	\item I dati inseriti dall'utente sono validati lato client (es. campi obbligatori, lunghezza minima password).
	\item I token di autenticazione sono gestiti in sessionStorage, riducendo i rischi di XSS.
	\item La comunicazione WebSocket è autenticata tramite token.
\end{itemize}

\section{Motivazioni progettuali}
La struttura a componenti e l'uso di Context rendono il frontend facilmente estendibile e manutenibile. Ogni funzione è isolata e testabile. L'adozione di WebSocket per la chat garantisce reattività e scalabilità. L'approccio "separation of concerns" tra presentazione, logica e stato migliora la qualità del codice e la collaborazione tra sviluppatori.